\subsection{Depth and connectivity}

Lemma 4.--- Let A be a Noetherian ring, F a closed subset of Spec(A), U the complementary open set, and u : M → N a homomorphism of finitely generated A-modules.

\begin{enumerate}
\def\labelenumi{\alph{enumi})}
\item
  Suppose that \(u_{\mathfrak{p}}: \mathrm{M}_{\mathfrak{p}} \to \mathrm{N}_{\mathfrak{p}}\) is injective for all \(\mathfrak{p} \in \mathrm{U}\) and suppose we have \(\operatorname{prof}_{\mathrm{F}}(\mathrm{M}) \geqslant 1\). Then \(u\) is injective.
\item
  Suppose that \(u_{\mathfrak{p}}: \mathrm{M}_{\mathfrak{p}} \to \mathrm{N}_{\mathfrak{p}}\) is bijective for all \(\mathfrak{p} \in \mathrm{U}\) and suppose we have \(\operatorname{prof}_{\mathrm{F}}(\mathrm{M}) \geqslant 2\) and \(\operatorname{prof}_{\mathrm{F}}(\mathrm{N}) \geqslant 1\). Then u is bijective.
\item
  The hypotheses imply Supp(Ker \(u\) ) \(\subset\) F, then \(\operatorname{Hom}_{\Lambda}(\operatorname{Ker} u, \mathbf{M}) = 0\) ( \(n^{\circ}\) 5, remark 1); we therefore have \(\operatorname{Ker} u = 0\) .
\item
  We already know that \(u\) is injective, and we have \(\operatorname{Supp}(\operatorname{Coker} u) \subset \mathbb{F}\) . By loc. cit., we have \(\operatorname{Hom}_{\mathbf{A}}(\operatorname{Coker} u, \mathbf{N}) = 0\) and \(\operatorname{Ext}_{\mathbf{A}}^{1}(\operatorname{Coker} u, \mathbf{M}) = 0\) . From the exact sequence of extension modules
\end{enumerate}

\[
\operatorname{Hom} _ {\mathrm{A}} (\operatorname{Coker} u, \mathrm{N}) \rightarrow \operatorname{Hom} _ {\mathrm{A}} (\operatorname{Coker} u, \operatorname{Coker} u) \rightarrow \operatorname{Ext} _ {\mathrm{A}} ^ {1} (\operatorname{Coker} u, \mathrm{M})
\]

we deduce \(\operatorname{Hom}_{\mathrm{A}}(\operatorname{Coker} u, \operatorname{Coker} u) = 0\) , which implies \(\operatorname{Coker} u = 0\) .

Remark 1.—Let A be a Noetherian ring, F a closed subset of \(\operatorname{Spec}(A)\) , U the complementary open set. For one to have \(\operatorname{prof}_{\mathrm{F}}(A) \geqslant 1\) , it is necessary and sufficient that one has \(\operatorname{Ass}(A) \subset U\) (remark 2, no. 5). When this condition is satisfied, each irreducible component of \(\operatorname{Spec}(A)\) meets U, so that U is dense in \(\operatorname{Spec}(A)\) .

THEOREM 3 (Hartshorne).--- Let A be a Noetherian ring, F a closed subset of Spec(A), and U the complementary open set. Suppose that one has prof \(_{F}\) (A) ≥ 2. Then, for every connected component Y of Spec(A), the set Y ∩ U is connected and dense in Y.

Suppose first that \(\operatorname{Spec}(A)\) is connected. By Remark 1, \(U\) is dense in \(\operatorname{Spec}(A)\), and it remains to prove that it is connected. Suppose, for a contradiction, that two disjoint nonempty open subsets \(U_0\) and \(U_1\) of \(\operatorname{Spec}(A)\) have union \(U\). Since \(\operatorname{Ass}(A)\) is contained in \(U\) by Remark 1, it is a disjoint union of \(\operatorname{Ass}(A) \cap U_0\) and \(\operatorname{Ass}(A) \cap U_1\). By IV, § 1, No. 1, Prop. 4, there exist ideals \(J_0\) and \(J_1\) of A such that \(\operatorname{Ass}(J_i) = \operatorname{Ass}(A) \cap U_i\), \(\operatorname{Ass}(A/J_i) = \operatorname{Ass}(A) \cap U_{1-i}\quad (i=0,1)\). The complement of \(U_i\) in \(\operatorname{Spec}(A)\) contains \(\operatorname{Ass}(A/J_i)\) and \(\operatorname{Ass}(J_{1-i})\); since it is closed, it contains \(\operatorname{Supp}(A/J_i)\) and \(\operatorname{Supp}(J_{1-i})\). For \(\mathfrak{p} \in U_i\), we thus have \((J_i)_{\mathfrak{p}} = A_{\mathfrak{p}}\) and \((J_{1-i})_{\mathfrak{p}} = 0\); this implies in particular that \(J_0\) and \(J_1\) are distinct from A. Let B be the A-module \(A/J_0 \times A/J_1\) and let \(u: A \to B\) be the canonical homomorphism. By what precedes, the homomorphism \(u_{\mathfrak{p}}\) is bijective for every \(\mathfrak{p} \in U\); moreover, we have \(\operatorname{Ass}(B) \subset U_0 \cup U_1 = U\), hence \(\operatorname{prof}_F(B) \geqslant 1\) by Remark 1. Lemma 4 then implies that \(u\) is bijective, which contradicts the connectedness of \(\operatorname{Spec}(A)\).

Let us treat the general case. Let J be an ideal of A such that \(F = V(J)\), and let Y be a connected component of \(\operatorname{Spec}(A)\). By II, § 4, No. 3, Prop. 15, there exists an idempotent element \(f\) of A such that Y is identified with the part \(\operatorname{Spec}(A_f)\) of \(\operatorname{Spec}(A)\). Then \(Y \cap F\) is identified with \(V(J_f)\); we have \(\operatorname{prof}_{A_f}(J_f, A_f) \geqslant \operatorname{prof}_A(J; A) \geqslant 2\) by Prop. 6, a) of No. 3. It follows from the first part of the proof that \(Y \cap U = Y - (Y \cap F)\) is connected and dense in Y.

COROLLARY 1.--- The map which associates to each connected component of U its closure in Spec(A) is a bijection from the set of connected components of U onto the set of connected components of Spec(A).

COROLLARY 2.--- For every Noetherian local ring B of depth \(\geqslant\) 2, the topological space \(\operatorname{Spec}(\mathbf{B}) = \{\mathfrak{m}_{\mathbf{B}}\}\) is connected.

COROLLARY 3.--- Under the hypotheses of Theorem 3, suppose that \(\text{Spec}(A_{\mathfrak{p}})\) Let it be irreducible (resp. that \(A_{p}\) (let it be integral) for all \(p \in U\) ; then \(\text{Spec}(A_{\mathfrak{p}})\) is irreducible (resp. \(A_{p}\) is integral) for all \(p \in \text{Spec}(A)\) .

Let \((\mathrm{Y}_{i})_{i\in\mathrm{I}}\) be the (finite) family of irreducible components of \(\operatorname{Spec}(A)\). Let \(\mathfrak{p}\in U\); since \(\operatorname{Spec}(A_{\mathfrak{p}})\) is irreducible, \(\mathfrak{p}\) contains only one minimal prime ideal of A and therefore belongs to only one of the \(Y_i\) (II, § 4, No. 3, Cor. 2 of Prop. 14). The subsets \(Y_i\cap U\) are closed subsets of U, disjoint, and nonempty since U is dense in \(\operatorname{Spec}(A)\) and irreducible by II, § 4, No. 1, Prop. 7; these are therefore the connected components of U. Their closures \(Y_i\) are the connected components of \(\operatorname{Spec}(A)\) (Cor. 1). This proves that the connected components of \(\operatorname{Spec}(A)\) are irreducible, hence that \(\operatorname{Spec}(A_{\mathfrak{p}})\) is irreducible for every \(\mathfrak{p}\) (No. 8).

Suppose that \(A_{\mathfrak{q}}\) is integral for all \(\mathfrak{q} \in U\). Let \(\mathfrak{p} \in \operatorname{Spec}(A)\). Since \(\operatorname{Spec}(A_{\mathfrak{p}})\) is irreducible, the nilradical of \(A_{\mathfrak{p}}\) is the unique minimal prime ideal of \(A_{\mathfrak{p}}\); it therefore belongs to \(\operatorname{Ass}(A_{\mathfrak{p}})\) (IV, § 1, No. 3, Cor. 1 of Prop. 7), and consequently is equal to \(\mathfrak{q}A_{\mathfrak{p}}\), where \(\mathfrak{q}\) is a prime ideal associated with A (IV, § 1, No. 2, Cor. of Prop. 5). We have \(\mathfrak{q} \in U\) (Remark 1) and \(\mathfrak{q}A_{\mathfrak{q}} \in \operatorname{Ass}(A_{\mathfrak{q}})\) (loc. cit.); since \(A_{\mathfrak{q}}\) is integral, \(\mathfrak{q}\) is zero, therefore \(A_{\mathfrak{p}}\) is an integral domain.

COROLLARY 4.--- Let \(A\) be a Noetherian ring whose spectrum is connected. Suppose there exists an integer \(d \geqslant 1\) such that \(\operatorname{prof}(A_{\mathfrak{p}}) \geqslant 2\) for every prime ideal \(\mathfrak{p}\) of A of height \textgreater{} d.

\begin{enumerate}
\def\labelenumi{\alph{enumi})}
\item
  For every closed subset Z of \(\operatorname{Spec}(\mathbf{A})\) of codimension \(>d\) the space \(\operatorname{Spec}(\mathbf{A}) - \mathbf{Z}\) is connected.
\item
  Let Y and Y' be irreducible components of Spec(A). There exists a sequence \(X_{1}, X_{2}, \ldots, X_{n}\) of irreducible components of Spec(A) such that one has \(X_{1} = Y\) , \(X_{n} = Y'\) and, for \(i = 1, \ldots, n - 1\)
\end{enumerate}

\[
\operatorname{codim} \left(\mathrm{X} _ {i} \cap \mathrm{X} _ {i + 1}, \operatorname{Spec} (\mathrm{A})\right) \leqslant d.
\]

Let \(Z \subset \operatorname{Spec}(A)\) a closed subset of codimension \(>d\) . For every \(\mathfrak{p} \in Z\) , one has \(\dim(A_{\mathfrak{p}}) > d\) , hence \(\operatorname{prof}(A_{\mathfrak{p}}) \geqslant 2\) , which implies that \(\operatorname{prof}_{Z}(A)\) is \(\geqslant 2\) (no. 5, prop. 8) and therefore that \(\operatorname{Spec}(A) - Z\) is connected (Thm. 3).

Let Z denote the union of the sets \(X^{\prime} \cap X^{\prime\prime}\) where \((X^{\prime}, X^{\prime\prime})\) traverses the (finite) set of pairs of irreducible components of Spec(A) such that \(\text{codim}(X^{\prime} \cap X^{\prime\prime}, \text{Spec}(A)) > d\) . By a), the set \(\text{Spec}(A) - Z\) is connected. All irreducible components of \(\text{Spec}(A)\) meet \(\text{Spec}(A) - Z\) ; their traces on \(\text{Spec}(A) - Z\) are the irreducible components of \(\text{Spec}(A) - Z\) (II, § 4, no. 1, prop. 7). Since \(\text{Spec}(A) - Z\) is connected, there exists a sequence \(X_{1}, \ldots, X_{n}\) of irreducible components of \(\text{Spec}(A)\) such that one has \(X_{1} - Z = Y - Z\) , \(X_{n} - Z = Y' - Z\) and \((X_{i} - Z) \cap (X_{i+1} - Z) \neq \emptyset\) for \(1 \leqslant i \leqslant n - 1\) ; in other words, we have \(X_{1} = Y\) , \(X_{n} = Y'\) and \(\text{codim}(X_{i} \cap X_{i+1}) \leqslant d\) .

Remark 2. When A is a Macaulay ring, one may apply the corollary with \(d=1\) ( \(\S\) 2, no. 5).

Example (Complete intersection formed by four coordinate planes of a 4-dimensional affine space).--- Let \(k\) be a field. Let S denote the polynomial ring \(k[\mathrm{T}_{1},\mathrm{T}_{2},\mathrm{T}_{3},\mathrm{T}_{4}]\). Recall (VIII, § 2, No. 4, Th. 3) that every maximal chain of prime ideals of S has length 4. Let \(\mathfrak{m}\) be the ideal of S generated by the \(\mathrm{T}_{i}\), \(\mathfrak{a}\) the ideal generated by \(\mathrm{T}_{1}\mathrm{T}_{2}\) and \(\mathrm{T}_{3}\mathrm{T}_{4}\), and, for \(1 \leqslant i < j \leqslant 4\), let \(\mathfrak{p}_{ij}\) be the ideal generated by \(\mathrm{T}_{i}\) and \(\mathrm{T}_{j}\). The ideals \(\mathfrak{p}_{ij}\) are prime of height 2, their sum is the maximal ideal \(\mathfrak{m}\) and one has \(\mathfrak{a} = \mathfrak{p}_{13} \cap \mathfrak{p}_{14} \cap \mathfrak{p}_{23} \cap \mathfrak{p}_{24}\).

\begin{enumerate}
\def\labelenumi{\alph{enumi})}
\item
  The ring A = S/a is reduced; the irreducible components of Spec(A) are the sets \(X_{ij} = V(\mathfrak{p}_{ij}/\mathfrak{a})\) for i = 1, 2, j = 3, 4, which are of dimension 2 and all contain the point m/a. In particular, Spec(A) is connected and of dimension 2. The intersection of two distinct components \(X_{ij}\) and \(X_{kl}\) is reduced to \(\{m/a\}\) if \(\{i,j\} \cap \{k,l\} = \varnothing\) , of dimension 1 otherwise. It follows that the conclusion of cor. 4 is satisfied with d = 1 (we shall see later that A is a Macaulay ring, so that the hypothesis of corollary 4 is itself satisfied for d = 1).
\item
  Let b be the ideal of S generated by \(T_{1}T_{2}\) , \(T_{1}T_{3}\) , \(T_{2}T_{4}\) , \(T_{3}T_{4}\) , and B the ring S/b. We have \(b = p_{14}p_{23} = p_{14} \cap p_{23}\) . The ring B is reduced. Its spectrum is identified with the closed subset \(X_{14} \cup X_{23}\) of Spec(A); it has two irreducible components (of dimension 2) whose intersection is reduced to a point. The depth of B along this point is strictly positive because B is reduced, and less than 1 by th. 3, hence equal to 1. Thus B is not a Macaulay ring.
\end{enumerate}

